\chapter{Plan de pruebas}
\label{anexo:plan_pruebas}

Este anexo define el plan de pruebas seguido para validar el funcionamiento de TierraViva. El objetivo es comprobar de forma progresiva cada parte del sistema antes de realizar una integración completa: sensores, comunicaciones, almacenamiento, dashboard y actuación sobre el riego.

Al tratarse de un prototipo que combina hardware, software y servicios externos, las pruebas se han organizado por fases. Esta estrategia reduce el riesgo de errores difíciles de aislar, ya que permite validar primero cada bloque por separado y después comprobar el comportamiento del sistema completo.

\section{Estrategia general de pruebas}

La validación de TierraViva se plantea como un proceso incremental. En primer lugar se prueban los componentes de forma aislada; después se comprueba la comunicación entre bloques; finalmente se valida el funcionamiento integrado del prototipo.

\begin{table}[H]
\centering
\renewcommand{\arraystretch}{1.25}
\small
\begin{tabular}{|c|p{0.30\textwidth}|p{0.43\textwidth}|}
\hline
\rowcolor{gray!20}
\textbf{Fase} & \textbf{Tipo de prueba} & \textbf{Objetivo principal} \\
\hline
F1 & Pruebas eléctricas básicas & Verificar alimentación, masas comunes, tensiones y conexiones antes de cargar el sistema completo. \\
\hline
F2 & Pruebas de sensores & Comprobar que los sensores entregan lecturas coherentes y estables. \\
\hline
F3 & Pruebas del módulo LTE & Validar comunicación AT, registro en red y envío de datos. \\
\hline
F4 & Pruebas de ThingSpeak & Confirmar publicación y lectura remota de telemetría. \\
\hline
F5 & Pruebas de Raspberry Pi & Verificar consulta de datos, almacenamiento SQLite, API y dashboard. \\
\hline
F6 & Pruebas de actuación sin carga real & Validar la lógica de riego usando relé o indicador, sin bomba conectada. \\
\hline
F7 & Pruebas de actuación con bomba & Comprobar activación temporizada y apagado seguro del sistema de riego. \\
\hline
F8 & Pruebas de integración completa & Validar el flujo completo estación--ThingSpeak--Raspberry--dashboard--riego. \\
\hline
F9 & Pruebas en entorno real o controlado & Evaluar el comportamiento del prototipo en una parcela, maceta o entorno agrícola simulado. \\
\hline
\end{tabular}
\caption{Estrategia general de pruebas del proyecto TierraViva.}
\label{tab:estrategia_pruebas}
\end{table}

\section{Entorno de pruebas}

El entorno de pruebas se compone de una estación remota, una Raspberry Pi como nodo de supervisión, una cuenta de ThingSpeak y un sistema de riego controlado mediante relé.

\begin{table}[H]
\centering
\renewcommand{\arraystretch}{1.25}
\small
\begin{tabular}{|p{0.28\textwidth}|p{0.58\textwidth}|}
\hline
\rowcolor{gray!20}
\textbf{Elemento} & \textbf{Uso durante las pruebas} \\
\hline
Arduino UNO & Ejecución del firmware de estación, lectura de sensores y control del relé. \\
\hline
A7670E-LASA & Comunicación LTE/4G mediante comandos AT y envío de telemetría. \\
\hline
BME280 & Medición de temperatura, humedad relativa y presión atmosférica. \\
\hline
BH1750 & Medición de luminosidad ambiental. \\
\hline
Sensor capacitivo de humedad del suelo & Medición de humedad del sustrato para la decisión de riego. \\
\hline
Módulo relé & Conmutación de la bomba o carga de prueba. \\
\hline
Bomba o electroválvula de 12 V & Actuación física sobre el sistema de riego. \\
\hline
ThingSpeak & Recepción de telemetría y consulta remota desde Raspberry Pi. \\
\hline
Raspberry Pi & Ejecución de scripts de consulta, almacenamiento SQLite, API y dashboard. \\
\hline
SQLite & Registro histórico de lecturas, estados y eventos de riego. \\
\hline
Dashboard web & Visualización del estado del sistema y de las últimas lecturas. \\
\hline
\end{tabular}
\caption{Elementos empleados en el entorno de pruebas.}
\label{tab:entorno_pruebas}
\end{table}

\section{Criterios de aceptación}

Para considerar válida una prueba, se definen criterios de aceptación claros. Estos criterios permiten evitar una validación subjetiva basada únicamente en que el sistema ``parece funcionar''.

\begin{table}[H]
\centering
\renewcommand{\arraystretch}{1.25}
\small
\begin{tabular}{|p{0.30\textwidth}|p{0.56\textwidth}|}
\hline
\rowcolor{gray!20}
\textbf{Área} & \textbf{Criterio de aceptación} \\
\hline
Alimentación & Las tensiones medidas se mantienen dentro de valores seguros antes y durante la prueba. \\
\hline
Sensores & Las lecturas son coherentes con el entorno y no presentan valores imposibles o fuera de rango. \\
\hline
Comunicación LTE & El módulo responde a comandos AT, registra red y permite enviar datos. \\
\hline
ThingSpeak & Los datos enviados por la estación aparecen en el canal correspondiente. \\
\hline
Raspberry Pi & El sistema consulta lecturas, las almacena localmente y mantiene ejecución estable. \\
\hline
API & Los endpoints devuelven respuestas válidas y coherentes con la base de datos. \\
\hline
Dashboard & La interfaz muestra datos actualizados y estados interpretables para el usuario. \\
\hline
Relé & El relé cambia de estado únicamente cuando corresponde y vuelve a estado seguro. \\
\hline
Bomba & La bomba se activa durante el tiempo previsto y se apaga al finalizar. \\
\hline
Seguridad & El sistema evita riegos repetidos, excesivos o provocados por datos no válidos. \\
\hline
\end{tabular}
\caption{Criterios generales de aceptación de pruebas.}
\label{tab:criterios_aceptacion}
\end{table}

\section{Casos de prueba}

La Tabla \ref{tab:casos_prueba} recoge los casos de prueba principales. Se han seleccionado pruebas representativas de cada bloque del sistema, priorizando aquellas que validan la funcionalidad esencial del prototipo.

\begin{table}[H]
\centering
\renewcommand{\arraystretch}{1.22}
\scriptsize
\begin{tabular}{|c|p{0.20\textwidth}|p{0.30\textwidth}|p{0.28\textwidth}|}
\hline
\rowcolor{gray!20}
\textbf{ID} & \textbf{Prueba} & \textbf{Procedimiento} & \textbf{Resultado esperado} \\
\hline
P01 & Verificación de alimentación & Medir la salida del LM2596 y comprobar masa común antes de conectar todos los módulos. & Tensiones correctas y ausencia de calentamiento o reinicios. \\
\hline
P02 & Lectura de BME280 & Ejecutar el firmware con el sensor conectado al bus I2C. & Temperatura, humedad y presión con valores coherentes. \\
\hline
P03 & Lectura de BH1750 & Exponer el sensor a diferentes condiciones de luz. & Variación de lux coherente entre sombra y luz directa. \\
\hline
P04 & Lectura de humedad del suelo & Probar el sensor en seco, húmedo y condición intermedia. & Valores diferenciables que permitan definir umbrales. \\
\hline
P05 & Comunicación AT con A7670E-LASA & Enviar comandos básicos como \texttt{AT}, \texttt{AT+CPIN?}, \texttt{AT+CSQ} y \texttt{AT+CEREG?}. & Respuestas válidas, SIM lista y registro en red. \\
\hline
P06 & Publicación en ThingSpeak & Enviar una lectura de prueba desde la estación. & La lectura aparece en el canal de la estación. \\
\hline
P07 & Consulta desde Raspberry Pi & Ejecutar el script de lectura de ThingSpeak. & La Raspberry recupera la última lectura disponible. \\
\hline
P08 & Almacenamiento en SQLite & Insertar lecturas obtenidas desde ThingSpeak en la base de datos local. & Las lecturas quedan registradas con estación, fecha y valores. \\
\hline
P09 & API local & Consultar endpoints como \texttt{/api/latest}, \texttt{/api/readings} y \texttt{/api/health}. & La API responde sin error y devuelve datos coherentes. \\
\hline
P10 & Dashboard web & Abrir la interfaz desde navegador y comprobar actualización de datos. & El dashboard muestra lecturas, estado y frescura de datos. \\
\hline
P11 & Relé sin bomba & Activar el pin de control y observar el cambio de estado del relé sin carga hidráulica. & El relé conmuta y vuelve a apagado al finalizar. \\
\hline
P12 & Bomba con activación temporizada & Conectar bomba o carga real y ejecutar una orden de riego limitada. & La bomba se activa durante el tiempo definido y se apaga correctamente. \\
\hline
P13 & Cooldown entre riegos & Forzar dos decisiones de riego consecutivas. & La segunda activación se bloquea si no ha pasado el tiempo mínimo. \\
\hline
P14 & Lectura inválida de humedad & Simular valor fuera de rango o desconexión del sensor. & El sistema no activa el riego y registra estado de error. \\
\hline
P15 & Pérdida temporal de comunicación & Interrumpir temporalmente la publicación o consulta de datos. & El sistema registra el fallo y mantiene estado seguro. \\
\hline
P16 & Integración completa & Ejecutar estación, ThingSpeak, Raspberry, API, dashboard y actuación. & El flujo completo funciona de extremo a extremo. \\
\hline
\end{tabular}
\caption{Casos de prueba principales definidos para TierraViva.}
\label{tab:casos_prueba}
\end{table}

\section{Secuencia de validación recomendada}

Para reducir riesgos durante las pruebas, la validación se realiza de forma gradual. No se conecta la bomba real hasta haber comprobado previamente la lógica de control y el comportamiento del relé.

\begin{table}[H]
\centering
\renewcommand{\arraystretch}{1.25}
\small
\begin{tabular}{|c|p{0.37\textwidth}|p{0.39\textwidth}|}
\hline
\rowcolor{gray!20}
\textbf{Paso} & \textbf{Acción} & \textbf{Motivo} \\
\hline
1 & Probar sensores y alimentación por separado. & Evitar que un error de cableado afecte al sistema completo. \\
\hline
2 & Probar el módulo LTE con comandos AT básicos. & Confirmar que la comunicación funciona antes de integrarla en el firmware. \\
\hline
3 & Publicar datos de prueba en ThingSpeak. & Validar el canal de telemetría antes de automatizar el envío. \\
\hline
4 & Consultar datos desde Raspberry Pi. & Comprobar que la supervisión remota recibe la información esperada. \\
\hline
5 & Validar base de datos, API y dashboard. & Confirmar que el backend local interpreta y muestra correctamente los datos. \\
\hline
6 & Activar relé sin bomba conectada. & Probar la lógica de actuación sin riesgo hidráulico. \\
\hline
7 & Activar bomba durante tiempos cortos y controlados. & Validar la actuación real con riesgo limitado. \\
\hline
8 & Ejecutar prueba completa con sensores, comunicación, backend y riego. & Confirmar el funcionamiento integrado del prototipo. \\
\hline
9 & Repetir la prueba en entorno real o controlado. & Comprobar comportamiento fuera del entorno de laboratorio. \\
\hline
\end{tabular}
\caption{Secuencia recomendada de validación del prototipo.}
\label{tab:secuencia_validacion}
\end{table}

\section{Pruebas de seguridad funcional}

La seguridad funcional en TierraViva se centra en evitar actuaciones no deseadas sobre el sistema de riego. No se trata de seguridad industrial certificada, sino de medidas básicas necesarias para un prototipo que acciona una bomba o electroválvula.

\begin{table}[H]
\centering
\renewcommand{\arraystretch}{1.25}
\small
\begin{tabular}{|p{0.35\textwidth}|p{0.45\textwidth}|p{0.12\textwidth}|}
\hline
\rowcolor{gray!20}
\textbf{Condición comprobada} & \textbf{Comportamiento esperado} & \textbf{Prioridad} \\
\hline
Arranque o reinicio de la estación & El relé debe comenzar apagado. & Alta \\
\hline
Orden o decisión de riego con duración excesiva & La duración debe limitarse al máximo permitido. & Alta \\
\hline
Lectura de humedad inválida & El sistema no debe activar la bomba. & Alta \\
\hline
Riegos consecutivos en poco tiempo & El sistema debe aplicar periodo de seguridad. & Alta \\
\hline
Error durante la ejecución & El relé debe apagarse de forma explícita. & Alta \\
\hline
Pérdida temporal de comunicación & La estación debe mantener un estado seguro. & Media \\
\hline
Duplicidad de orden o evento & La estación no debe repetir una actuación ya ejecutada. & Media \\
\hline
\end{tabular}
\caption{Pruebas de seguridad funcional del sistema de riego.}
\label{tab:pruebas_seguridad_funcional}
\end{table}

\section{Relación entre riesgos y pruebas}

El plan de pruebas se ha definido teniendo en cuenta los riesgos principales del proyecto. La Tabla \ref{tab:trazabilidad_riesgos_pruebas} muestra la relación entre algunos riesgos identificados y las pruebas previstas.

\begin{table}[H]
\centering
\renewcommand{\arraystretch}{1.25}
\small
\begin{tabular}{|c|p{0.42\textwidth}|p{0.34\textwidth}|}
\hline
\rowcolor{gray!20}
\textbf{Riesgo} & \textbf{Descripción} & \textbf{Pruebas asociadas} \\
\hline
R01 & Alimentación insuficiente o inestable del módulo LTE. & P01, P05, P06 \\
\hline
R02 & Lecturas incorrectas del sensor de humedad del suelo. & P04, P13, P14 \\
\hline
R03 & Degradación puntual de la comunicación LTE. & P05, P06, P15 \\
\hline
R04 & Indisponibilidad temporal de ThingSpeak. & P06, P07, P15 \\
\hline
R05 & Activación incorrecta o excesiva de la bomba. & P11, P12, P13 \\
\hline
R06 & Daño eléctrico por cableado o integración incorrecta. & P01, P11, P16 \\
\hline
R07 & Exposición ambiental del montaje. & P16 y prueba en entorno controlado \\
\hline
R08 & Exposición accidental de claves o credenciales. & Revisión de repositorio y ficheros de configuración \\
\hline
R09 & Inconsistencias entre firmware, backend, API y dashboard. & P07, P08, P09, P10, P16 \\
\hline
\end{tabular}
\caption{Trazabilidad entre riesgos y pruebas.}
\label{tab:trazabilidad_riesgos_pruebas}
\end{table}

\section{Registro de resultados}

Los resultados de cada prueba deben registrarse para facilitar la trazabilidad del proyecto. Para cada ensayo se recomienda anotar al menos:

\begin{itemize}
    \item identificador de la prueba;
    \item fecha de ejecución;
    \item versión del firmware o software utilizado;
    \item configuración relevante;
    \item resultado obtenido;
    \item incidencias detectadas;
    \item acciones correctivas aplicadas;
    \item resultado de la repetición, si procede.
\end{itemize}

Este registro se desarrolla de forma complementaria en el anexo dedicado al diario de experimentos, donde se documentan las pruebas realizadas, los problemas encontrados y las modificaciones aplicadas durante el desarrollo.

\section{Estados de una prueba}

Para homogeneizar el registro de resultados se utilizarán los siguientes estados:

\begin{table}[H]
\centering
\renewcommand{\arraystretch}{1.25}
\small
\begin{tabular}{|p{0.22\textwidth}|p{0.64\textwidth}|}
\hline
\rowcolor{gray!20}
\textbf{Estado} & \textbf{Significado} \\
\hline
No ejecutada & La prueba está definida pero aún no se ha realizado. \\
\hline
Superada & La prueba cumple todos los criterios de aceptación. \\
\hline
Superada con observaciones & La funcionalidad principal funciona, pero se han detectado aspectos a mejorar o documentar. \\
\hline
Fallida & La prueba no cumple el resultado esperado y requiere corrección. \\
\hline
Bloqueada & La prueba no puede ejecutarse por falta de componente, entorno, configuración o dependencia previa. \\
\hline
\end{tabular}
\caption{Estados definidos para el seguimiento de pruebas.}
\label{tab:estados_prueba}
\end{table}

\section{Conclusión}

El plan de pruebas de TierraViva permite validar el sistema de forma progresiva y controlada. La estrategia planteada evita conectar directamente todos los componentes sin haber comprobado antes alimentación, sensores, comunicación, almacenamiento, API, dashboard y actuación.

Además, la relación entre riesgos y pruebas refuerza la calidad del proyecto, ya que cada punto crítico identificado dispone de una comprobación asociada. De este modo, la validación del prototipo no se limita a demostrar que funciona en condiciones ideales, sino que también contempla errores previsibles, situaciones de fallo y medidas de seguridad funcional.